\section{26 --- The Reliability Decision Framework}

Reliability engineering is not the optimization of a single scalar score. It is a disciplined, multi-dimensional decision process that balances correctness, tail latency, cost, and provenance under uncertainty:

\begin{figure}[h!]
\centering
\begin{tikzpicture}[
  scale=0.92, every node/.style={transform shape},
  box/.style={draw=primary, fill=white, rounded corners=3pt, font=\footnotesize\sffamily, align=center, inner sep=5pt, line width=0.8pt},
  flow/.style={-{Latex[length=2.5mm]}, line width=1.2pt, color=secondary}
]
\node[box, fill=blue!5, text width=9.5cm] (ev) at (0, 0) {\textbf{EMPIRICAL EVIDENCE}\\Paired Seed Executions $\cdot$ Linked Traces};

\node[box, fill=teal!5, text width=5.6cm, below left=0.6cm and -0.8cm of ev] (corr) {\textbf{Correctness Gate}\\$\bullet$ Grounded Pass Rate\\$\bullet$ Citation Completeness\\$\bullet$ Zero Hallucinated Tokens};

\node[box, fill=warning!5, text width=5.6cm, below right=0.6cm and -0.8cm of ev] (ops) {\textbf{Operational Gate}\\$\bullet$ p95 / p99 Latency Budget\\$\bullet$ Cost per Resolved Query\\$\bullet$ Attempt Amplification Ceiling};

\node[box, fill=purple!5, text width=9.5cm, below=0.6cm of $(corr.south)!0.5!(ops.south)$] (unc) {\textbf{Uncertainty Quantification}\\$\bullet$ Wilson 95\% Score Bounds $\cdot$ 10,000-Resample Bootstrap CIs $\cdot$ McNemar Significance};

\node[box, fill=secondary!10, text width=9.5cm, below=0.5cm of unc] (rep) {\textbf{Reproducibility \& Provenance}\\$\bullet$ Byte-for-Byte Replay Stability $\cdot$ Fail-Closed Claim Audit};

\node[box, draw=success, fill=success!15, line width=1.2pt, text width=9.5cm, below=0.5cm of rep] (dec) {\textbf{\color{success!80!black}RELIABILITY DECISION}\\Eligible for Production Deployment $\cdot$ Bounded Envelope};

\draw[flow] (ev.south) -| (corr.north);
\draw[flow] (ev.south) -| (ops.north);
\draw[flow] (corr.south) |- (unc.west);
\draw[flow] (ops.south) |- (unc.east);
\draw[flow] (unc.south) -- (rep.north);
\draw[flow] (rep.south) -- (dec.north);
\end{tikzpicture}
\caption{The Multi-Dimensional Reliability Decision Framework.}
\end{figure}

\subsection{The Five Decision Gating Rules}
\begin{enumerate}
    \item \textbf{Rule 1 (Correctness Floor):} The Wilson 95\% confidence interval lower bound for grounded pass rate must exceed the operational threshold ($w_{\text{lower}} \ge 0.78$). A policy with high point success but wide uncertainty is rejected.
    \item \textbf{Rule 2 (Safety / Wrong-Answer Ceiling):} The Wilson 95\% upper bound on wrong user-visible answers must remain below the risk tolerance ($w_{\text{upper}} \le 0.02$). Returning a wrong answer is strictly penalized over abstaining.
    \item \textbf{Rule 3 (Tail Latency Ceiling):} The bootstrap 95\% upper bound on p95 latency must remain within the user SLA ($L_{p95} \le 80\text{ units}$).
    \item \textbf{Rule 4 (Efficiency Pareto Dominance):} Among eligible policies passing Rules 1--3, candidate Policy A is preferred over Policy B if its paired bootstrap differences in success-per-dollar and success-per-second strictly exclude zero.
    \item \textbf{Rule 5 (Stress Envelope Certification):} The selected policy must undergo adversarial stress injection (worse fault severity, correlated outages, tight budgets). If it fails Rules 1--3 under stress, its production deployment must be gated behind dynamic envelope monitoring.
\end{enumerate}

% ==============================================================================
\vspace{12pt}
\section{27 --- What FAULTLINE Taught Me}

These lessons are not motivational aphorisms; they are hard-won empirical truths derived from debugging distributed AI agent systems:

\begin{enumerate}
    \item \textbf{Dashboards Hide Semantic Failure.} High infrastructure availability is easily achieved by returning empty or degraded answers. An AI observability system that does not evaluate the semantic grounding of the output is a compliance facade.
    
    \item \textbf{Deterministic Environments Make Science Possible.} You cannot debug stochastic agent failures in a stochastic environment. By locking the world, documents, and tool contracts with seeds, every observed failure becomes an attributable causal event.
    
    \item \textbf{Recovery is Not Synonymous with Reliability.} Adding retries, fallbacks, and circuit breakers often reduces true system reliability by causing retry storms, latency inflation, and the silent delivery of degraded answers.
    
    \item \textbf{Evaluators Cannot Validate Themselves.} Using an uncalibrated LLM judge to measure agent hallucinations is circular. Evaluators must be benchmarked against independent human ground truth, and must be quarantined when they exhibit blind spots.
    
    \item \textbf{Averages Conceal Catastrophes.} An aggregate accuracy of 85\% can conceal a 0\% success rate on deep multi-hop queries or specific fault families. Reliability engineering requires micro-slicing data across fault types, hop depths, and severities.
    
    \item \textbf{Correlation Breaks Independent Architecture.} In a distributed system, assuming component independence during an outage is fatal. When upstream providers degrade, retries multiply load rather than buying success.
    
    \item \textbf{Reproducibility is Part of the Result, Not Documentation.} A scientific claim without an executable, containerized reproduction path is merely an anecdote. Research integrity requires proving that a stranger can clone the code and produce the exact same numbers.
    
    \item \textbf{Define Failure Before Measuring Success.} If you cannot enumerate the physical failure modes of your agent's tools, memory, and model boundaries, you cannot measure its reliability. System design begins with a failure taxonomy.
\end{enumerate}
